\documentclass[a4paper,11pt]{amsart}
\usepackage[margin=1in]{geometry}
\usepackage{amsmath,amssymb,amsthm}

\theoremstyle{plain}
\newtheorem{theorem}{Theorem}
\newtheorem{lemma}[theorem]{Lemma}
\newtheorem{proposition}[theorem]{Proposition}

\theoremstyle{definition}
\newtheorem{definition}[theorem]{Definition}
\newtheorem{remark}[theorem]{Remark}

\title{A Sidon set of size 12 in $[1,100]$}

\begin{document}
\maketitle

\begin{abstract}
A \emph{Sidon set} in a set of integers is one whose pairwise sums are all
distinct. We exhibit an explicit Sidon set of cardinality $12$ contained in
the interval $[1,100]$, namely
$\{1,2,4,8,17,28,40,59,69,77,94,99\}$. Validity is reduced via a short
lemma to the claim that all $\binom{12}{2}=66$ positive pairwise
differences are distinct, which is verified by displaying the entire
difference table and observing that the resulting $66$ integers in
$[1,98]$ contain no repetition. Consequently, the maximum cardinality of
a Sidon set in $[1,100]$ is at least $12$, improving on the value $11$
previously recorded for this instance.
\end{abstract}

\section{Setting and result}

\begin{definition}
For an integer $N\geq 1$ let $[1,N]:=\{1,2,\dots,N\}$. A finite set
$S\subseteq\mathbb{Z}$ is a \emph{Sidon set} (or $B_2$\nobreakdash-set) if
\[
a+b=c+d\quad\text{with}\quad a,b,c,d\in S\quad\Longrightarrow\quad
\{a,b\}=\{c,d\}.
\]
Equivalently, the map $\{a,b\}\mapsto a+b$ from two\nobreakdash-element
subsets of $S$ to $\mathbb{Z}$ is injective. Define
\[
F(N):=\max\{|S| : S\subseteq[1,N]\text{ is a Sidon set}\}.
\]
\end{definition}

The aim of this note is to prove the following.

\begin{theorem}\label{thm:main}
$F(100)\geq 12$.
\end{theorem}

\section{The construction}

Let
\begin{equation}\label{eq:S}
S:=\{1,\,2,\,4,\,8,\,17,\,28,\,40,\,59,\,69,\,77,\,94,\,99\}.
\end{equation}
We list its elements in increasing order as $a_1<a_2<\cdots<a_{12}$, that
is
\[
(a_1,a_2,\dots,a_{12})=(1,2,4,8,17,28,40,59,69,77,94,99).
\]

\begin{proposition}\label{prop:size}
$|S|=12$ and $S\subseteq[1,100]$.
\end{proposition}

\begin{proof}
The display \eqref{eq:S} consists of $12$ strictly increasing positive
integers, so $S$ has exactly $12$ pairwise distinct elements. Its
smallest element is $1\geq 1$ and its largest is $99\leq 100$, hence
$S\subseteq[1,100]$.
\end{proof}

\section{Reduction to differences}

We use the following standard equivalence between distinctness of sums
and distinctness of positive differences.

\begin{lemma}\label{lem:diff}
Let $S\subseteq\mathbb{Z}$ be a finite set. Then $S$ is a Sidon set if
and only if, for any two two\nobreakdash-element subsets
$\{a,b\},\{c,d\}\subseteq S$ with $a<b$ and $c<d$,
\[
\{a,b\}\neq\{c,d\}\quad\Longrightarrow\quad b-a\neq d-c.
\]
\end{lemma}

\begin{proof}
We prove the contrapositive of each implication.

\smallskip
\noindent\emph{(Sidon $\Rightarrow$ differences distinct.)}
Suppose $\{a,b\}\neq\{c,d\}$ are two\nobreakdash-element subsets of $S$
with $a<b$, $c<d$, and $b-a=d-c$. Then
\begin{equation}\label{eq:starsum}
a+d=b+c.
\end{equation}
If $a=c$ then $b=d$, contradicting $\{a,b\}\neq\{c,d\}$, so $a\neq c$.
Renaming the two pairs if necessary, we may assume $a<c$. Combined
with $b-a=d-c$ this gives $b=a+(d-c)<c+(d-c)=d$, hence $b<d$.
Therefore
\[
a<c,\qquad a<b,\qquad b<d,\qquad c<d,
\]
so $a$ is strictly smaller than each of $b,c,d$, and $d$ is strictly
larger than each of $a,b,c$. The two pairs $\{a,d\}$ and $\{b,c\}$
are therefore disjoint, and in particular distinct. By
\eqref{eq:starsum} they have the same sum, contradicting the Sidon
property.

\smallskip
\noindent\emph{(Differences distinct $\Rightarrow$ Sidon.)}
Suppose $\{a,b\}\neq\{c,d\}$ are two\nobreakdash-element subsets of $S$
with $a<b$, $c<d$, and $a+b=c+d$. If $a=c$ then $b=d$, contradiction;
hence $a\neq c$, and after renaming we may assume $a<c$. Subtracting
from $a+b=c+d$ gives $b-d=c-a>0$, so $b>d$. Combined with $c<d$ we
obtain
\[
a<c<d<b,
\]
and the pairs $\{a,c\}$ and $\{d,b\}$ are disjoint and hence distinct.
They realise the common positive difference
\[
c-a=b-d,
\]
contradicting the assumption that the positive differences of $S$ are
pairwise distinct.
\end{proof}

\section{Verification of the Sidon property}

By Lemma~\ref{lem:diff}, to show that the set $S$ of \eqref{eq:S} is
Sidon it suffices to verify that the $\binom{12}{2}=66$ positive
differences
\[
d_{ij}:=a_j-a_i,\qquad 1\leq i<j\leq 12,
\]
are pairwise distinct. The full difference table follows; the entry in
row~$i$, column~$j$ is $d_{ij}=a_j-a_i$, and each entry is verified by
direct subtraction from
$(a_1,\dots,a_{12})=(1,2,4,8,17,28,40,59,69,77,94,99)$.

\medskip
\begin{center}
\renewcommand{\arraystretch}{1.05}
\begin{tabular}{c|rrrrrrrrrrr}
$i\,\backslash\,j$ & $2$ & $3$ & $4$ & $5$ & $6$ & $7$ & $8$ & $9$ & $10$ & $11$ & $12$\\
\hline
$1$  & $1$ & $3$ & $7$ & $16$ & $27$ & $39$ & $58$ & $68$ & $76$ & $93$ & $98$\\
$2$  &     & $2$ & $6$ & $15$ & $26$ & $38$ & $57$ & $67$ & $75$ & $92$ & $97$\\
$3$  &     &     & $4$ & $13$ & $24$ & $36$ & $55$ & $65$ & $73$ & $90$ & $95$\\
$4$  &     &     &     & $9$  & $20$ & $32$ & $51$ & $61$ & $69$ & $86$ & $91$\\
$5$  &     &     &     &      & $11$ & $23$ & $42$ & $52$ & $60$ & $77$ & $82$\\
$6$  &     &     &     &      &      & $12$ & $31$ & $41$ & $49$ & $66$ & $71$\\
$7$  &     &     &     &      &      &      & $19$ & $29$ & $37$ & $54$ & $59$\\
$8$  &     &     &     &      &      &      &      & $10$ & $18$ & $35$ & $40$\\
$9$  &     &     &     &      &      &      &      &      & $8$  & $25$ & $30$\\
$10$ &     &     &     &      &      &      &      &      &      & $17$ & $22$\\
$11$ &     &     &     &      &      &      &      &      &      &      & $5$
\end{tabular}
\end{center}
\medskip

Every entry is a positive integer in $[1,98]$; indeed
$d_{ij}\geq d_{j-1,j}\geq 1$ and $d_{ij}\leq d_{1,12}=a_{12}-a_1=98$.
Sorting the $66$ tabulated values in increasing order yields
\begin{equation}\label{eq:diffs-sorted}
\begin{aligned}
&1,\,2,\,3,\,4,\,5,\,6,\,7,\,8,\,9,\,10,\,11,\,12,\,13,\,15,\,16,\,17,\,18,\,19,\,20,\,22,\\
&23,\,24,\,25,\,26,\,27,\,29,\,30,\,31,\,32,\,35,\,36,\,37,\,38,\,39,\,40,\,41,\,42,\\
&49,\,51,\,52,\,54,\,55,\,57,\,58,\,59,\,60,\,61,\,65,\,66,\,67,\,68,\,69,\,71,\,73,\\
&75,\,76,\,77,\,82,\,86,\,90,\,91,\,92,\,93,\,95,\,97,\,98.
\end{aligned}
\end{equation}
The sequence \eqref{eq:diffs-sorted} is strictly increasing and
contains $66$ distinct integers (one may verify the cardinality by
counting, for example as $13+6+6+4+8+1+2+2+5+5+1+1+3+1+1+4+1+2=66$).
Therefore the positive differences of $S$ are pairwise distinct.

A few features of the table merit a comment, since they could
superficially appear suspicious. The element $40=a_7$ also appears as
a difference: $40=99-59=a_{12}-a_8$, so $a_7+a_8=a_{12}$. Likewise
$a_4+a_9=8+69=77=a_{10}$ and $a_5+a_{10}=17+77=94=a_{11}$. The Sidon
condition forbids only the coincidence of two distinct two\nobreakdash-element
subset sums; it does \emph{not} forbid an element of $S$ from being the
sum of two other elements of $S$, nor does it forbid an element of $S$
from equalling a difference of two other elements. Each of the
relations $a_7+a_8=a_{12}$, $a_4+a_9=a_{10}$, $a_5+a_{10}=a_{11}$
involves three distinct elements of $S$, and so corresponds to a
\emph{single} pairwise sum (namely $a_7+a_8$, $a_4+a_9$, or
$a_5+a_{10}$ respectively). It would only contradict Sidonicity if the
same value were also expressible as a different pairwise sum
$a_i+a_j$ with $\{i,j\}$ disjoint from the original pair, and
\eqref{eq:diffs-sorted} (via Lemma~\ref{lem:diff}) rules this out.

By Lemma~\ref{lem:diff}, $S$ is a Sidon set.

\section{Conclusion}

\begin{proof}[Proof of Theorem~\ref{thm:main}]
Proposition~\ref{prop:size} gives $S\subseteq[1,100]$ with $|S|=12$, and
the verification of Section~4 shows that $S$ is a Sidon set. Hence
$F(100)\geq|S|=12$.
\end{proof}

The bound $F(100)\geq 12$ improves on the previously recorded
$F(100)\geq 11$ by one. The argument is entirely elementary: it consists
of the explicit witness $S$, the equivalence between sum\nobreakdash-collisions
and difference\nobreakdash-collisions (Lemma~\ref{lem:diff}), and the
inspection of the $66$\nobreakdash-entry difference table summarised in
\eqref{eq:diffs-sorted}.

\begin{remark}
Equivalently, one could verify the Sidon property of $S$ directly from
the definition by listing the $66$ pairwise sums $a_i+a_j$
($1\leq i<j\leq 12$), all of which lie in $[3,193]$, and checking that
they are pairwise distinct. The two checks are equivalent by
Lemma~\ref{lem:diff}; the difference formulation is presented because
its entries are smaller, fit comfortably within $[1,98]$, and are
therefore more readily inspected.
\end{remark}

\end{document}
